% SPDX-License-Identifier: CC-BY-SA-4.0
% Author: Matthieu Perrin

\begingroup

\begin{exercice}[Algorithme de Cocke-Younger-Kasami]\label{exo:parsing/decision/CYK}

  Soient les grammaires algébriques suivantes~:
  $$G_1\eqdef\left\langle \{a, b\}, \{S, X, A, B\}, S,
    \left\{\begin{array}{lll}
    S   &\rightarrow& AB \mid  AX\\
    X   &\rightarrow& SB \\
    A   &\rightarrow& a \\
    B   &\rightarrow& b
    \end{array}\right\}\right\rangle,$$
    $$G_2\eqdef\left\langle\{a, b\}, \{S, X, A, B\}, S,
    \left\{\begin{array}{lllllll}
    S   &\rightarrow& AX \mid  AS\\
    X   &\rightarrow& b \mid  XB\\
    A   &\rightarrow& a \\
    B   &\rightarrow& b
    \end{array}\right\}\right\rangle.$$

  \begin{question}
  \item Pour chacune de ces grammaires, trouver 2 mots appartenant au langage et 2 mots n'y appartenant pas.
  \item Utiliser l'algorithme CYK pour déterminer si $aabb$ et $aaabb$ sont générés par ces grammaires.
  \item Déterminer les langages engendrés par ces grammaires.
  \end{question}

\end{exercice}

\begin{correction}

  \begin{itemize}
  \item Pour $G_1$, le langage est $\{a^nb^n \mid  n > 0\}$ (de type 2).
    L'analyse CYK de aabb:
    $$\begin{array}{|c|c|c|c|c|}
      \hline
      & T[...,1] & T[...,2] & T[...,3] & T[...,4]\\
      \hline
      T[1,...] & A & A & B & B \\
      \cline{1-5}
      T[2,...] & \emptyset  & S  & \emptyset \\
      \cline{1-4}
      T[3,...] & \emptyset & X  \\
      \cline{1-3}
      T[4,...]   & S \\
      \cline{1-2}
    \end{array}$$
    $aabb$ appartient au langage.
    L'analyse CYK de aaabb:
    $$\begin{array}{|c|c|c|c|c|c|}
      \hline
      & T[...,1] & T[...,2] & T[...,3] & T[...,4] & T[...,5]\\
      \hline
      T[1,...] & A & A & A & B & B \\
      \cline{1-6}
      T[2,...] & \emptyset  & \emptyset  & S  & \emptyset \\
      \cline{1-5}
      T[3,...] & \emptyset  & \emptyset & X  \\
      \cline{1-4}
      T[4,...]   & \emptyset  & S \\
      \cline{1-3}
      T[5,...]   & \emptyset  \\
      \cline{1-2}
    \end{array}$$
    $aaabb$ n'appartient pas au langage.

  \item
    Pour $G_2$, le langage est $\{a^nb^m \mid  n > 0, m > 0\}$  (de type 3)
    L'analyse CYK de aabb:
    $$\begin{array}{|c|c|c|c|c|}
      \hline
      & T[...,1] & T[...,2] & T[...,3] & T[...,4]\\
      \hline
      T[1,...] & A & A & B,X & B,X \\
      \cline{1-5}
      T[2,...] & \emptyset  & S  & X \\
      \cline{1-4}
      T[3,...] & S & S  \\
      \cline{1-3}
      T[4,...]   & S \\
      \cline{1-2}
    \end{array}$$
    $aabb$ appartient au langage.
    L'analyse CYK de aaabb:
    $$\begin{array}{|c|c|c|c|c|c|}
      \hline
      & T[...,1] & T[...,2] & T[...,3] & T[...,4] & T[...,5]\\
      \hline
      T[1,...] & A & A & A & B,X & B,X \\
      \cline{1-6}
      T[2,...] & \emptyset  & \emptyset  & S  & X \\
      \cline{1-5}
      T[3,...] & \emptyset  & S & S  \\
      \cline{1-4}
      T[4,...]   & S  & S \\
      \cline{1-3}
      T[5,...]   & S  \\
      \cline{1-2}
    \end{array}$$
    $aaabb$ appartient au langage.
  \end{itemize}

\end{correction}

\endgroup
\endinput
